\documentclass[border=10pt,tikz]{standalone}
\usepackage{tikz}
\usetikzlibrary{positioning,arrows.meta,shapes,fit,backgrounds}
\usepackage{kotex}

% 디버거가 Cortex-M의 halt·step·breakpoint·watchpoint 자원에 닿는 경로.
% 출처: OpenOCD src/target/adi_v5_swd.c / adi_v5_jtag.c (DP 접근),
%       src/target/armv7m.h (debug_ap = MEM-AP), src/target/cortex_m.c
%       (DHCSR 읽기·쓰기는 mem_ap_*_u32(debug_ap, DCB_DHCSR ...);
%        FP_CTRL에서 fp_num_code, DWT_CTRL에서 dwt_num_comp를 읽음),
%       src/target/cortex_m.h (C_STEP = BIT(2), C_MASKINTS = BIT(3)).

\begin{document}
% _design.tex — shared TikZ design system for all blog diagrams
% Inputs nothing on its own — meant to be \input{} from each diagram .tex
% AFTER \begin{document}, BEFORE \begin{tikzpicture}.
%
% Provides a unified visual language across GoF / DSA / EMC++ / Beautiful series:
%   - Colors: muted, accessible, color-coded by role
%   - Node styles: consistent sizing, rounded, subtle borders
%   - Edge styles: UML-style inheritance, aggregation, plus dashed alternates
%
% Usage:
%   \documentclass[border=10pt,tikz]{standalone}
%   \usepackage{tikz}
%   \usetikzlibrary{positioning,arrows.meta,shapes,fit,backgrounds}
%   \begin{document}
%   \input{../_design.tex}     % adjust relative path
%   \begin{tikzpicture}[blog]  % apply 'blog' style
%   ...

\definecolor{absfill}{RGB}{220,232,247}   % cool blue
\definecolor{absbord}{RGB}{72,118,182}
\definecolor{confill}{RGB}{249,222,222}   % warm red
\definecolor{conbord}{RGB}{197,93,93}
\definecolor{impfill}{RGB}{223,238,222}   % soft green
\definecolor{impbord}{RGB}{99,154,93}
\definecolor{clifill}{RGB}{251,239,205}   % pale yellow
\definecolor{clibord}{RGB}{200,167,72}
\definecolor{hifill} {RGB}{253,224,200}   % accent orange
\definecolor{hibord} {RGB}{210,135,68}
\definecolor{nodefill}{RGB}{250,250,250}
\definecolor{nodebord}{RGB}{120,120,120}
\definecolor{edgec}  {RGB}{80,80,80}
\definecolor{edgesc} {RGB}{145,145,145}
\definecolor{groupbord}{RGB}{200,200,200}

\tikzset{
  blog/.style={
    font=\sffamily\small,
    every node/.style={font=\sffamily\small}
  },
  % --- node roles ---
  cls/.style={
    rectangle, rounded corners=2pt,
    draw=nodebord, line width=0.4pt,
    minimum height=8mm, minimum width=28mm,
    inner sep=3pt, fill=nodefill, align=center
  },
  abs/.style={cls, draw=absbord, fill=absfill, font=\sffamily\itshape\small},
  con/.style={cls, draw=conbord, fill=confill},
  imp/.style={cls, draw=impbord, fill=impfill},
  cli/.style={cls, draw=clibord, fill=clifill},
  hi/.style={cls, draw=hibord, fill=hifill, font=\sffamily\bfseries\small},
  % --- circle (for trees / graph nodes) ---
  vx/.style={
    circle, draw=absbord, line width=0.4pt,
    minimum size=8mm, fill=absfill, inner sep=0pt
  },
  vxc/.style={vx, draw=conbord, fill=confill},
  vxh/.style={vx, draw=hibord, fill=hifill, font=\sffamily\bfseries},
  % --- decision diamond ---
  q/.style={
    diamond, draw=clibord, line width=0.4pt, fill=clifill,
    inner sep=2pt, aspect=2.4, align=center
  },
  % --- edges ---
  inh/.style={                                  % UML inheritance
    -{Triangle[length=3mm,width=3mm,open]},
    draw=edgec, line width=0.5pt
  },
  agg/.style={                                  % UML aggregation
    {Diamond[length=2.5mm,width=2.5mm,open]}-{Stealth[length=2.2mm]},
    draw=edgec, line width=0.45pt
  },
  uses/.style={                                 % directed reference
    -{Stealth[length=2.2mm]}, draw=edgec, line width=0.5pt
  },
  arr/.style={                                  % plain arrow (graphs/trees)
    -{Stealth[length=2mm]}, draw=edgec, line width=0.45pt
  },
  edge/.style={                                 % undirected (trees)
    draw=edgec, line width=0.45pt
  },
  alt/.style={                                  % dashed: similar/alt
    -{Stealth[length=2mm]}, dashed, draw=edgesc, line width=0.45pt
  },
  ret/.style={                                  % dashed creates/returns
    -{Stealth[length=2mm]}, dashed, draw=edgesc, line width=0.45pt
  },
  thr/.style={                                  % threading link (red dashed)
    ->, dashed, draw=conbord, line width=0.5pt
  },
  % --- group / region ---
  group/.style={
    rectangle, rounded corners=4pt,
    draw=groupbord, line width=0.4pt,
    inner sep=10pt, fill=black!2
  },
  glabel/.style={font=\sffamily\bfseries\small, anchor=north west, inner sep=2pt}
}

% Korean font convenience macro — included by every Hangul diagram
\providecommand{\KOR}{}


\begin{tikzpicture}[blog,
    layer/.style={rectangle, rounded corners=3pt, draw=absbord, fill=absfill,
                  minimum width=70mm, minimum height=12mm, align=center,
                  font=\sffamily\small},
    probe/.style={layer, draw=hibord, fill=hifill},
    block/.style={rectangle, rounded corners=2pt, draw=impbord, fill=impfill,
                  text width=40mm, minimum height=22mm, align=center,
                  font=\sffamily\small},
    arr/.style={-{Stealth[length=2.4mm]}, draw=edgec, line width=0.55pt},
    lbl/.style={font=\sffamily\footnotesize, text=inkc, fill=white, inner sep=2pt}
]

\node[probe] (probe) at (0, 6.6) {\textbf{probe + 서버} (OpenOCD · J-Link GDB Server)};

\node[layer] (dp) at (0, 4.4)
  {\textbf{DP} (Debug Port)\\[2pt]\footnotesize SWD 또는 JTAG로 접근};
\node[layer] (ap) at (0, 2.4)
  {\textbf{MEM-AP} (Access Port)\\[2pt]\footnotesize 디버그 레지스터를 메모리 주소로 읽고 씀};

\draw[arr] (probe) -- node[lbl, right=1mm]{SWD 또는 JTAG} (dp);
\draw[arr] (dp) -- (ap);

\node[block] (dcb) at (-4.8, -0.6)
  {\textbf{core debug (DHCSR)}\\[2pt]\footnotesize C\_HALT = bit 1\\[2pt]\footnotesize C\_STEP = bit 2\\[2pt]\footnotesize C\_MASKINTS = bit 3};
\node[block] (fpb) at ( 0,   -0.6)
  {\textbf{FPB}\\[2pt]\footnotesize hardware breakpoint 비교기\\[2pt]\footnotesize 개수는 FP\_CTRL에서 읽음};
\node[block] (dwt) at ( 4.8, -0.6)
  {\textbf{DWT}\\[2pt]\footnotesize watchpoint 비교기\\[2pt]\footnotesize 개수는 DWT\_CTRL에서 읽음};

\coordinate (bus) at (0, 1.0);
\draw[arr] (ap.south) -- (bus) -| (dcb.north);
\draw[arr] (bus) -- (fpb.north);
\draw[arr] (bus) -| (dwt.north);

\node[font=\sffamily\footnotesize, text=mutec, align=center] at (0, -2.6)
  {비교기 수는 core·구현마다 다르므로 서버가 위 레지스터를 읽어 확인합니다.};

\end{tikzpicture}
\end{document}
